% =====================================================================
% Capítulo 3. Trabajos relacionados                Presupuesto: 6 páginas
% ---------------------------------------------------------------------
% No es una lista de resúmenes. Cada trabajo se cuenta con 3 preguntas:
%   ¿qué hicieron? ¿qué obtuvieron (con número)? ¿en qué difiere de esta tesis?
% Agrupa por tema, no por autor. Termina con una tabla comparativa y un
% párrafo que nombre la BRECHA que cubre la tesis.
% Usa \cite{clave} siempre; nunca cites algo que no hayas leído.
% =====================================================================
\chapter{Trabajos relacionados}
\label{cap:trabajos}

% Introducción del capítulo (3--4 líneas): qué grupos de trabajos se revisan y por qué.

\section{Seguridad de LoRaWAN y ataques de \textit{bit flipping}}
\label{sec:rel-lorawan}
% alizadeh2023bitflipping (ICM): qué datos, cómo generan los ataques (bits
% aleatorios), qué reportan. Diferencia: no relacionan la posición del bit con
% el daño en la decisión. + otros análisis de seguridad de LoRaWAN 1.1.

\section{Detección de anomalías en redes de monitoreo atmosférico}
\label{sec:rel-calidad-aire}
% allka2024paraad (PARAAD) y allka2024stcraad (STC-RAAD): mismo grupo, O3 en
% Barcelona. Qué anomalías sintéticas usan (sesgo, deriva), qué métricas.
% PARAAD es el comparador de esta tesis: decirlo aquí.

\section{Control de calidad con contexto espacial}
\label{sec:rel-espacial}
% lussana2010sct (SCT), baserud2020titan (TITAN), dereszynski2011spatiotemporal,
% deng2021gdn (GDN): cómo usan a los vecinos y qué pasa cuando un vecino
% también está mal (motiva el filtro de NERAD).

\section{Detección de intrusiones en IoT}
\label{sec:rel-ids}
% garciasaez2025fedwla (FedWLA): ataca el tráfico de red, no el dato; la
% prevalencia de ataques (28--49 %) no es comparable con 0.77 %.

\section{Síntesis y brecha}
\label{sec:rel-brecha}
% Tabla \ref{tab:comparacion-trabajos}: filas = trabajos; columnas = datos,
% supervisado, usa vecinas, amenaza, resultado reportado.
% Párrafo final: qué combinación no existe y esta tesis aporta.
